\chapter*{ВВЕДЕНИЕ}
\addcontentsline{toc}{chapter}{ВВЕДЕНИЕ}

Настоящий документ представляет собой промежуточный научно-технический отчёт о
выполнении первого этапа научно-исследовательской и опытно-конструкторской
работы (НИОКР) по теме <<Разработка и тестирование прототипа программного
комплекса для моделирования и симуляции роботов с использованием графических
ускорителей>>, выполняемой при поддержке Фонда содействия инновациям по
программе <<Старт-1>>. Отчёт оформлен в соответствии с требованиями
ГОСТ~7.32~\cite{gost732}.

\textbf{Актуальность работы.}
Робототехника является одной из наиболее динамично развивающихся областей
современной техники. Проектирование, верификация и тестирование
робототехнических систем требуют значительных временны\'{х} и финансовых затрат
при работе с физическим оборудованием. Программные комплексы моделирования и
симуляции позволяют перенести значительную часть этих работ в виртуальную среду,
сокращая цикл разработки, снижая риски повреждения дорогостоящего оборудования и
обеспечивая воспроизводимость экспериментов. Современные сценарии применения ---
обучение с подкреплением, построение цифровых двойников, отработка алгоритмов
управления --- предъявляют всё более высокие требования к производительности
симуляции, что делает применение графических ускорителей (GPU) для расчёта
физической динамики не только желательным, но и необходимым.

В условиях ограничений на использование зарубежного программного обеспечения
актуальность создания отечественного программного комплекса робототехнической
симуляции существенно возрастает. Разрабатываемый комплекс ориентирован на
импортонезависимую инженерную разработку и объединяет средства загрузки и
визуализации моделей роботов и сцен, автоматического построения кинематических
схем, физической симуляции с GPU-ускорением и графический интерфейс
пользователя.

\textbf{Цель этапа.}
Целью первого этапа является разработка и тестирование отдельных подсистем
прототипа программного комплекса: модулей загрузки и отрисовки моделей роботов и
сцен, алгоритма построения кинематической схемы, алгоритмов симуляции на
центральном (CPU) и графическом (GPU) процессорах, а также базового графического
интерфейса пользователя, с проведением предварительного тестирования на
искусственных данных и сбором эталонных показателей производительности.

\textbf{Задачи этапа} соответствуют работам календарного плана:
\begin{enumerate}
  \item разработка и реализация алгоритма загрузки и отрисовки моделей роботов в
        формате URDF;
  \item реализация модуля загрузки моделей сцен и объектов окружения в форматах
        SDF и STL;
  \item разработка и программная реализация алгоритма генерации кинематической
        схемы на основе формата STEP;
  \item разработка и реализация алгоритма симуляции, не использующего
        графический процессор;
  \item разработка и тестирование базового графического интерфейса пользователя;
  \item предварительное тестирование прототипа программного комплекса на
        искусственных данных и сбор эталонных показателей;
  \item разработка и программная реализация алгоритма численного интегрирования
        дифференциальных уравнений с помощью GPU.
\end{enumerate}

\textbf{Место этапа в выполнении НИОКР в целом.}
Первый этап закладывает функциональный фундамент программного комплекса. На нём
создаются независимо тестируемые подсистемы (загрузка моделей, построение
кинематики, физическое ядро, интерфейс) и формируется набор эталонных
показателей, относительно которых на последующих этапах будет оцениваться
прирост производительности и точности при масштабировании GPU-ускорения и
интеграции подсистем в единый прикладной комплекс.

\textbf{Структура отчёта.}
Основная часть отчёта построена в соответствии с календарным планом. Раздел~1
обосновывает выбор направления исследований и общую методику. Разделы~2--8
описывают результаты по каждой из работ этапа. Раздел~9 содержит обобщение и
оценку полученных результатов. В заключении приведены краткие выводы и оценка
полноты решения поставленных задач.

\textbf{Программные результаты этапа} размещены в следующих репозиториях:
\begin{itemize}
  \item \texttt{modeuler} --- загрузка и отрисовка моделей в формате URDF,
        загрузка и отрисовка сцен в форматах SDF и STL;
  \item \texttt{kinechain} --- алгоритм генерации кинематической схемы из файлов
        в формате STEP;
  \item \texttt{modeuler-win} --- базовый графический интерфейс пользователя;
  \item \texttt{dyphur} --- алгоритм симуляции и решатель (солвер)
        дифференциальных уравнений с помощью GPU.
\end{itemize}
